% ==============================================================================
% חלק מקדים: תקציר מנהלים, תוכן עניינים ורשימות (Front Matter)
% ==============================================================================

\chapter*{תקציר מנהלים}
\addcontentsline{toc}{chapter}{תקציר מנהלים}

דוח מכין זה מציג את התשתית התיאורטית, המחקרית וההנדסית של פלטפורמת \projectname{} -- מערכת מחקר עתידית מבוססת סוכני בינה מלאכותית אוטונומיים (\textenglish{Multi-Agent System - MAS}), המיועדת לחקור, למדל ולשקף למשתמשים צעירים את השפעותיה של צריכת תוכן אלגוריתמי ברשתות חברתיות על תהליכי השוואה חברתית, דימוי עצמי ורווחה נפשית.

במהלך העשור האחרון, הפכו פלטפורמות מבוססות המלצה (דוגמת אינסטגרם וטיקטוק) לזירה המרכזית שבה מתעצב עולמם החברתי והרגשי של מתבגרים וצעירים. לצד היתרונות התקשורתיים, ספרות אמפירית ענפה מצביעה על מתאם מדאיג בין חשיפה רציפה לתכנים מלוטשים ואידיאליים לבין עלייה במדדי דיכאון, חרדה, חוסר שביעות רצון מדימוי הגוף (\textenglish{Body Image Dissatisfaction}) ושחיקה בערך העצמי. הבעיה היסודית אינה נובעת אך ורק ממשך זמן השהייה במסך, אלא מ\textbf{אטימות קוגניטיבית} (\textenglish{Cognitive Opacity}): המשתמש חווה תהליכי השוואה חברתית כלפי מעלה (\textenglish{Upward Social Comparison}) באופן רציף, מבלי שיהיה מודע למנגנוני האופטימיזציה המסחרית ולשיווק המשפיענים המעצבים את תמונת המציאות המוצגת בפניו.

כלי הרווחה הדיגיטלית הקיימים כיום (כגון \textenglish{Apple Screen Time} או הגבלות זמן באפליקציות) מציעים מענה גס ופטרנליסטי המתמקד במדידה כמותית בלבד, ללא התייחסות לאופי התוכן, להקשר האינטראקציה או למצבו הנפשי של המשתמש. \projectname{} מציעה שינוי פרדיגמה יסודי: מעבר מהגבלת זמן פסיבית ל"מראה קוגניטיבית" מעצימה. המערכת מבוססת על חלוקת תפקידים בין ארבעה סוכני \textenglish{LLM} אוטונומיים: סוכן תצפית (\agentObserver{}), סוכן ניתוח (\agentAnalyst{}), סוכן רפלקציה (\agentReflector{}) וסוכן התערבות (\agentInterventionist{}). סוכנים אלו מנטרים את גירויי ההשוואה, ממדלים את הדינמיקה הרגשית ומספקים פיגומים מטא-קוגניטיביים המאפשרים למשתמש לפתח מודעות ביקורתית וחוסן נפשי.

דוח מכין זה פורס את התשתית התיאורטית והמחקרית של הפרויקט: החל ממיפוי ההקשר החברתי-טכנולוגי והגדרת בעיית האטימות הקוגניטיבית, דרך סקירה ביקורתית של תיאוריית ההשוואה החברתית של פסטינגר ותיאוריית הפער העצמי של היגינס, ועד לניתוח הספרות האמפירית, חשיפת הכשל של כלי רווחה קיימים והתוויית הפרדיגמה המושגית של מערכת מרובת הסוכנים.

\vspace{0.4cm}

\begin{conceptbox}{חזון פלטפורמת \projectname{}}
מעבר ממדידה כמותית פטרנליסטית של "שעות מסך" אל \textbf{"מראה קוגניטיבית משקפת ומעצימה"}:
פלטפורמת מחקר עתידית מבוססת סוכנים אוטונומיים המפרקת את האטימות האלגוריתמית, מבודדת גירויי השוואה ושיווק משפיענים, ומספקת פיגומים מטא-קוגניטיביים המאפשרים לצעירים לפתח חוסן נפשי ובחירה דיגיטלית מודעת.
\end{conceptbox}

\newpage

\section*{מילון מונחים וסימונים מרכזיים}
\addcontentsline{toc}{section}{מילון מונחים וסימונים מרכזיים}

\begin{table}[H]
\centering
\small
\begin{tabularx}{\textwidth}{>{\bfseries\color{NavyBlueDark}}p{3.4cm} p{5.0cm} X}
\toprule
\textbf{מונח בעברית} & \textbf{מונח לועזי} & \textbf{הגדרה ומשמעות במערכת \projectname{}} \\
\midrule
מערכת מרובת סוכנים & \textenglish{Multi-Agent System (MAS)} & ארכיטקטורה חישובית של ישויות תוכנה אוטונומיות הפועלות בתיאום להשגת מטרות מורכבות. \\
השוואה חברתית כלפי מעלה & \textenglish{Upward Social Comparison} & תהליך פסיכולוגי שבו הפרט משווה את עצמו לאחרים הנתפסים כמוצלחים, מושכים או עשירים ממנו. \\
אפקט הניגוד & \textenglish{Contrast Effect} & תגובה רגשית שלילית שבה השוואה לאידיאל בלתי מושג מרחיבה את הפער העצמי הנתפס. \\
פער העצמי & \textenglish{Self-Discrepancy} & הפער הקוגניטיבי בין ה"עצמי הממשי" לבין ה"עצמי האידיאלי" (\textenglish{Actual vs. Ideal Self}). \\
מדד בולטות השוואה & \textenglish{$\comparisonScore$} & ציון כמותי מנורמל המחושב על ידי \agentAnalyst{} למדידת פוטנציאל ההשוואה בתוכן נתון. \\
מדד חשיפה למשפיענים & \textenglish{$\exposureIndex$} & מדד משוקלל המודד את נפח התוכן הממומן ואידיאליזציית סגנון החיים שנצרך ברצף זמן. \\
וקטור מצב רווחה נפשית & \textenglish{$\wellbeingVector$} & וקטור רב-ממדי דינמי המייצג את מצב הערך העצמי, המאזן הרגשי והעייפות הדיגיטלית של המשתמש. \\
סולם הערך העצמי של רוזנברג & \textenglish{RSES} & כלי פסיכומטרי מתוקף למדידת הערכה עצמית גלובלית ורגעית. \\
לוח רגשות חיוביים ושליליים & \textenglish{PANAS} & שאלון סטנדרטי להערכת מצבי רגש חיוביים ושליליים קצרי-טווח. \\
עיבוד מקומי בקצה & \textenglish{Local-First / Edge AI} & פרדיגמת מחשוב שבה כל ניתוח התוכן ומידול המצב מתבצעים על גבי המכשיר ללא זליגת מידע פרטי. \\
התנגדות פסיכולוגית & \textenglish{Psychological Reactance} & תגובה התנהגותית מרדנית של המשתמש בעת חסימה פטרנליסטית שרירותית של ממשק המשתמש. \\
\bottomrule
\end{tabularx}
\end{table}

\newpage
% תוכן עניינים ורשימות
{
  \hypersetup{linkcolor=NavyBlueDark}
  \tableofcontents
}

\newpage
{
  \hypersetup{linkcolor=NavyBlueDark}
  \listoffigures
  \vspace{0.8cm}
  \listoftables
}
